% 贯穿图上的行走、路径、闭行走与环。边旁数字是经过该边的步序。
\begin{tikzpicture}[>=Stealth,
  v/.style={circle,draw=FigureStroke,fill=FigureBlueFill,line width=.8pt,minimum size=6.5mm,inner sep=0pt},
  e/.style={draw=FigureGrid,line width=.8pt},
  h/.style={draw=FigureBlue,line width=1.6pt},
  s/.style={fill=white,inner sep=.8pt,text=FigureBlue,font=\figurefont\fontsize{7.5}{9}\selectfont}]
\newcommand{\base}{\node[v] (1) at (-1.0,0) {$1$}; \node[v] (2) at (0,.8) {$2$};
  \node[v] (3) at (0,-.8) {$3$}; \node[v] (4) at (1.0,0) {$4$};}
\begin{scope}
  \base
  \draw[e] (2)--(4) (3)--(4);
  \draw[h] (1)--node[s]{$1,4$}(2); \draw[h] (2)--node[s]{$2$}(3); \draw[h] (3)--node[s]{$3$}(1);
  \node[align=center] at (0,-1.65) {(a) 行走$1,2,3,1,2$\\长度$4$};
\end{scope}
\begin{scope}[xshift=2.75cm]
  \base
  \draw[e] (1)--(3) (2)--(3) (3)--(4);
  \draw[h] (1)--node[s]{$1$}(2); \draw[h] (2)--node[s]{$2$}(4);
  \node[align=center] at (0,-1.65) {(b) 路径$1,2,4$\\长度$2$};
\end{scope}
\begin{scope}[xshift=5.5cm]
  \base
  \draw[e] (1)--(2) (1)--(3) (2)--(3) (3)--(4);
  \draw[h] (2)--node[s]{$1,2$}(4);
  \node[align=center] at (0,-1.65) {(c) 闭行走$2,4,2$\\不是环};
\end{scope}
\begin{scope}[xshift=8.25cm]
  \base
  \draw[e] (2)--(4) (3)--(4);
  \draw[h] (1)--node[s]{$1$}(2); \draw[h] (2)--node[s]{$2$}(3); \draw[h] (3)--node[s]{$3$}(1);
  \node[align=center] at (0,-1.65) {(d) 环$1,2,3,1$\\长度$3$};
\end{scope}
\end{tikzpicture}
